%!TeX root=../gardentop.tex
\chapter{<Might I Have a Bit of Earth?>} 
	
\begin{figure}[t!]
\centering
\includegraphics[width=\linewidth]{february}
\end{figure}

 \lettrine[]{M}{ary} ran so fast that she was rather out of breath when she reached her room. Her hair was ruffled on her forehead and her cheeks were bright pink. Her dinner was waiting on the table, and Martha was waiting near it.

\zz
<Tha's a bit late,> she said. <Where has tha' been?>

\zz
<I've seen Dickon!> said Mary. <I've seen Dickon!>

<I knew he'd come,> said Martha exultantly. <How does tha' like him?>

<I think—I think he's beautiful!> said Mary in a determined voice.

Martha looked rather taken aback but she looked pleased, too.

<Well,> she said, <he's th' best lad as ever was born, but us never thought he was handsome. His nose turns up too much.>

<I like it to turn up,> said Mary.

<An' his eyes is so round,> said Martha, a trifle doubtful. <Though they're a nice colour.>

<I like them round,> said Mary. <And they are exactly the colour of the sky over the moor.>

Martha beamed with satisfaction.

<Mother says he made 'em that colour with always lookin' up at th' birds an' th' clouds. But he has got a big mouth, hasn't he, now?>

<I love his big mouth,> said Mary obstinately. <I wish mine were just like it.>

Martha chuckled delightedly.

<It'd look rare an' funny in thy bit of a face,> she said. <But I knowed it would be that way when tha' saw him. How did tha' like th' seeds an' th' garden tools?>

<How did you know he brought them?> asked Mary.

<Eh! I never thought of him not bringin' 'em. He'd be sure to bring 'em if they was in Yorkshire. He's such a trusty lad.>

Mary was afraid that she might begin to ask difficult questions, but she did not. She was very much interested in the seeds and gardening tools, and there was only one moment when Mary was frightened. This was when she began to ask where the flowers were to be planted.

<Who did tha' ask about it?> she inquired.

<I haven't asked anybody yet,> said Mary, hesitating.

<Well, I wouldn't ask th' head gardener. He's too grand, Mr~Roach is.>

<I've never seen him,> said Mary. <I've only seen under-gardeners and Ben Weatherstaff.>

<If I was you, I'd ask Ben Weatherstaff,> advised Martha. <He's not half as bad as he looks, for all he's so crabbed. Mr~Craven lets him do what he likes because he was here when Mrs~Craven was alive, an' he used to make her laugh. She liked him. Perhaps he'd find you a corner somewhere out o' the way.>

<If it was out of the way and no one wanted it, no one could mind my having it, could they?> Mary said anxiously.

<There wouldn't be no reason,> answered Martha. <You wouldn't do no harm.>

Mary ate her dinner as quickly as she could and when she rose from the table she was going to run to her room to put on her hat again, but Martha stopped her.

<I've got somethin' to tell you,> she said. <I thought I'd let you eat your dinner first. Mr~Craven came back this mornin' and I think he wants to see you.>

Mary turned quite pale.

<Oh!> she said. <Why! Why! He didn't want to see me when I came. I heard Pitcher say he didn't.>

<Well,> explained Martha, <Mrs~Medlock says it's because o' mother. She was walkin' to Thwaite village an' she met him. She'd never spoke to him before, but Mrs~Craven had been to our cottage two or three times. He'd forgot, but mother hadn't an' she made bold to stop him. I don't know what she said to him about you but she said somethin' as put him in th' mind to see you before he goes away again, to-morrow.>

<Oh!> cried Mary, <is he going away to-morrow? I am so glad!>

<He's goin' for a long time. He mayn't come back till autumn or winter. He's goin' to travel in foreign places. He's always doin' it.>

<Oh! I'm so glad—so glad!> said Mary thankfully.

If he did not come back until winter, or even autumn, there would be time to watch the secret garden come alive. Even if he found out then and took it away from her she would have had that much at least.

<When do you think he will want to see—>

She did not finish the sentence, because the door opened, and Mrs~Medlock walked in. She had on her best black dress and cap, and her collar was fastened with a large brooch with a picture of a man's face on it. It was a coloured photograph of Mr~Medlock who had died years ago, and she always wore it when she was dressed up. She looked nervous and excited.

<Your hair's rough,> she said quickly. <Go and brush it. Martha, help her to slip on her best dress. Mr~Craven sent me to bring her to him in his study.>

All the pink left Mary's cheeks. Her heart began to thump and she felt herself changing into a stiff, plain, silent child again. She did not even answer Mrs~Medlock, but turned and walked into her bedroom, followed by Martha. She said nothing while her dress was changed, and her hair brushed, and after she was quite tidy she followed Mrs~Medlock down the corridors, in silence. What was there for her to say? She was obliged to go and see Mr~Craven and he would not like her, and she would not like him. She knew what he would think of her.

She was taken to a part of the house she had not been into before. At last Mrs~Medlock knocked at a door, and when some one said, <Come in,> they entered the room together. A man was sitting in an armchair before the fire, and Mrs~Medlock spoke to him.

<This is Miss Mary, sir,> she said.

<You can go and leave her here. I will ring for you when I want you to take her away,> said Mr~Craven.

When she went out and closed the door, Mary could only stand waiting, a plain little thing, twisting her thin hands together. She could see that the man in the chair was not so much a hunchback as a man with high, rather crooked shoulders, and he had black hair streaked with white. He turned his head over his high shoulders and spoke to her.

<Come here!> he said.

Mary went to him.

He was not ugly. His face would have been handsome if it had not been so miserable. He looked as if the sight of her worried and fretted him and as if he did not know what in the world to do with her.

<Are you well?> he asked.

<Yes,> answered Mary.

<Do they take good care of you?>

<Yes.>

He rubbed his forehead fretfully as he looked her over.

<You are very thin,> he said.

<I am getting fatter,> Mary answered in what she knew was her stiffest way.

What an unhappy face he had! His black eyes seemed as if they scarcely saw her, as if they were seeing something else, and he could hardly keep his thoughts upon her.

<I forgot you,> he said. <How could I remember you? I intended to send you a governess or a nurse, or some one of that sort, but I forgot.>

<Please,> began Mary. <Please—> and then the lump in her throat choked her.

<What do you want to say?> he inquired.

<I am—I am too big for a nurse,> said Mary. <And please—please don't make me have a governess yet.>

He rubbed his forehead again and stared at her.

<That was what the Sowerby woman said,> he muttered absent-mindedly.

Then Mary gathered a scrap of courage.

<Is she—is she Martha's mother?> she stammered.

<Yes, I think so,> he replied.

<She knows about children,> said Mary. <She has twelve. She knows.>

He seemed to rouse himself.

<What do you want to do?>

<I want to play out of doors,> Mary answered, hoping that her voice did not tremble. <I never liked it in India. It makes me hungry here, and I am getting fatter.>

He was watching her.

<Mrs~Sowerby said it would do you good. Perhaps it will,> he said. <She thought you had better get stronger before you had a governess.>

<It makes me feel strong when I play and the wind comes over the moor,> argued Mary.

<Where do you play?> he asked next.

<Everywhere,> gasped Mary. <Martha's mother sent me a skipping-rope. I skip and run—and I look about to see if things are beginning to stick up out of the earth. I don't do any harm.>

<Don't look so frightened,> he said in a worried voice. <You could not do any harm, a child like you! You may do what you like.>

Mary put her hand up to her throat because she was afraid he might see the excited lump which she felt jump into it. She came a step nearer to him.

<May I\@?> she said tremulously.

Her anxious little face seemed to worry him more than ever.

<Don't look so frightened,> he exclaimed. <Of course you may. I am your guardian, though I am a poor one for any child. I cannot give you time or attention. I am too ill, and wretched and distracted; but I wish you to be happy and comfortable. I don't know anything about children, but Mrs~Medlock is to see that you have all you need. I sent for you to-day because Mrs~Sowerby said I ought to see you. Her daughter had talked about you. She thought you needed fresh air and freedom and running about.>

<She knows all about children,> Mary said again in spite of herself.

<She ought to,> said Mr~Craven. <I thought her rather bold to stop me on the moor, but she said—Mrs~Craven had been kind to her.> It seemed hard for him to speak his dead wife's name. <She is a respectable woman. Now I have seen you I think she said sensible things. Play out of doors as much as you like. It's a big place and you may go where you like and amuse yourself as you like. Is there anything you want?> as if a sudden thought had struck him. <Do you want toys, books, dolls?>

<Might I,> quavered Mary, <might I have a bit of earth?>

In her eagerness she did not realize how queer the words would sound and that they were not the ones she had meant to say. Mr~Craven looked quite startled.

<Earth!> he repeated. <What do you mean?>

<To plant seeds in—to make things grow—to see them come alive,> Mary faltered.

He gazed at her a moment and then passed his hand quickly over his eyes.

<Do you—care about gardens so much,> he said slowly.

<I didn't know about them in India,> said Mary. <I was always ill and tired and it was too hot. I sometimes made little beds in the sand and stuck flowers in them. But here it is different.>

Mr~Craven got up and began to walk slowly across the room.

<A bit of earth,> he said to himself, and Mary thought that somehow she must have reminded him of something. When he stopped and spoke to her his dark eyes looked almost soft and kind.

<You can have as much earth as you want,> he said. <You remind me of some one else who loved the earth and things that grow. When you see a bit of earth you want,> with something like a smile, <take it, child, and make it come alive.>

<May I take it from anywhere—if it's not wanted?>

<Anywhere,> he answered. <There! You must go now, I am tired.> He touched the bell to call Mrs~Medlock. <Good-by. I shall be away all summer.>

Mrs~Medlock came so quickly that Mary thought she must have been waiting in the corridor.

<Mrs~Medlock,> Mr~Craven said to her, <now I have seen the child I understand what Mrs~Sowerby meant. She must be less delicate before she begins lessons. Give her simple, healthy food. Let her run wild in the garden. Don't look after her too much. She needs liberty and fresh air and romping about. Mrs~Sowerby is to come and see her now and then and she may sometimes go to the cottage.>

Mrs~Medlock looked pleased. She was relieved to hear that she need not <look after> Mary too much. She had felt her a tiresome charge and had indeed seen as little of her as she dared. In addition to this she was fond of Martha's mother.

<Thank you, sir,> she said. <Susan Sowerby and me went to school together and she's as sensible and good-hearted a woman as you'd find in a day's walk. I never had any children myself and she's had twelve, and there never was healthier or better ones. Miss Mary can get no harm from them. I'd always take Susan Sowerby's advice about children myself. She's what you might call healthy-minded—if you understand me.>

<I understand,> Mr~Craven answered. <Take Miss Mary away now and send Pitcher to me.>

When Mrs~Medlock left her at the end of her own corridor Mary flew back to her room. She found Martha waiting there. Martha had, in fact, hurried back after she had removed the dinner service.

<I can have my garden!> cried Mary. <I may have it where I like! I am not going to have a governess for a long time! Your mother is coming to see me and I may go to your cottage! He says a little girl like me could not do any harm and I may do what I like—anywhere!>

<Eh!> said Martha delightedly, <that was nice of him wasn't it?>

<Martha,> said Mary solemnly, <he is really a nice man, only his face is so miserable and his forehead is all drawn together.>

She ran as quickly as she could to the garden. She had been away so much longer than she had thought she should and she knew Dickon would have to set out early on his five-mile walk. When she slipped through the door under the ivy, she saw he was not working where she had left him. The gardening tools were laid together under a tree. She ran to them, looking all round the place, but there was no Dickon to be seen. He had gone away and the secret garden was empty—except for the robin who had just flown across the wall and sat on a standard rose-bush watching her.

<He's gone,> she said wofully. <Oh! was he—was he—was he only a wood fairy?>

Something white fastened to the standard rose-bush caught her eye. It was a piece of paper—in fact, it was a piece of the letter she had printed for Martha to send to Dickon. It was fastened on the bush with a long thorn, and in a minute she knew Dickon had left it there. There were some roughly printed letters on it and a sort of picture. At first she could not tell what it was. Then she saw it was meant for a nest with a bird sitting on it. Underneath were the printed letters and they said:

<I will cum bak.>